\documentclass[11pt,a4paper]{article}

\usepackage[utf8]{inputenc}
\usepackage[T1]{fontenc}
\usepackage{amsmath,amssymb}
\usepackage{graphicx}
\usepackage{booktabs}
\usepackage[margin=2.5cm]{geometry}
\usepackage{siunitx}
\usepackage[colorlinks=true,linkcolor=blue,citecolor=blue,urlcolor=blue]{hyperref}

\newcommand{\csi}{CsI}
\newcommand{\bgo}{BGO}
\newcommand{\xzero}{X_0}

\title{\textbf{CryspLight}: GPU (Metal) simulation of scintillation-light transport
in wrapped \csi{} and \bgo{} crystals read out by a SiPM matrix}
\author{J.J.~G\'omez-Cadenas \and Claude (Anthropic)}
\date{July 13, 2026 \\[2mm] \small Version 0.2 --- design note}

\begin{document}
\maketitle

\begin{abstract}
We formulate the problem of simulating the generation and propagation of scintillation
photons in a monolithic \csi{} or \bgo{} crystal of transverse size
$48\times48~\mathrm{mm}^2$ and depth $2\xzero$, wrapped in a specular (3M Vikuiti ESR)
and/or diffuse (PTFE, ``Teflon'') reflector, and read out on the single unwrapped face by an
$8\times8$ matrix of $6\times6~\mathrm{mm}^2$ SiPMs. The gamma transport --- a 511~keV photon
interacting in the crystal via Compton scattering and photoelectric absorption, leaving a set
of localized energy deposits --- is solved by the existing PTCrysp engine, which this project
vendors. The new element is the \emph{optical} Monte Carlo: each deposit converts into
$\mathcal{O}(10^3\text{--}10^5)$ scintillation photons, emitted isotropically with an
exponential time profile ($\tau \simeq 800$~ns for \csi{} at 77~K, $\simeq 1500$~ns for \bgo{}
at 195~K), which then bounce inside the crystal until they are either absorbed (in the bulk or
at a wall) or detected at the SiPM plane. Because the photons are mutually independent, the
transport is embarrassingly parallel; we propose to implement it as a Metal compute kernel on
Apple-silicon GPUs, written in Julia via Metal.jl with a CPU-identical reference path. This
note collects the optical properties needed (crystal refractive index, bulk attenuation,
emission spectrum and time constants; reflector behaviour and how Geant4 parameterizes it),
specifies the surface physics model, defines the algorithm and its GPU mapping, estimates the
expected throughput, and lays out a validation plan against Geant4 and analytic limits.
The photon-detection efficiency $\mathrm{PDE}(\lambda)$ is applied in the kernel, so the
output is directly the detected light: per-SiPM photoelectron counts in 100~ns time bins plus
the first-photoelectron arrival time per SiPM. The baseline configuration for the initial
runs wraps both crystals in thick (1--2~mm) PTFE; the ESR study follows later.
\end{abstract}

\tableofcontents

%=====================================================================
\section{Problem statement}
\label{sec:problem}

\subsection{Setup}

A monolithic scintillating crystal, \csi{} (pure, operated cold, $\sim$77--100~K) or \bgo{}
(operated at $\sim$195~K, dry ice), has transverse dimensions $48\times48~\mathrm{mm}^2$ and
depth equal to two radiation lengths:
\begin{equation}
  L_z = 2\xzero =
  \begin{cases}
    37.2~\mathrm{mm} & \csi{} \quad (\xzero = 18.6~\mathrm{mm}),\\
    22.4~\mathrm{mm} & \bgo{} \quad (\xzero = 11.2~\mathrm{mm}).
  \end{cases}
\end{equation}
Five of the six faces (the front face, where the gamma enters, and the four lateral faces) are
wrapped in a reflector: 3M Vikuiti Enhanced Specular Reflector (ESR) film, PTFE (Teflon) tape,
or a combination. \textbf{The baseline for the initial runs is thick (1--2~mm) PTFE on both
crystals}; the ESR study is deferred (Sec.~\ref{sec:reflectors}). The wrapping is \emph{not}
optically coupled: a thin air gap separates the
polished crystal surface from the reflector, which has important consequences
(Sec.~\ref{sec:geant4}). The sixth face (the back face, $z = L_z$) is optically coupled ---
through a thin layer of optical grease, $n_g = 1.45$ --- to an $8\times8$
matrix of SiPMs, each with a $6\times6~\mathrm{mm}^2$ active area (e.g.\ Hamamatsu
S14160-6050HS~\cite{hamamatsu}). The array pitch of 6~mm exactly tiles the
$48\times48~\mathrm{mm}^2$ face; dead borders between devices are handled by a fill-factor
parameter.

\subsection{What is already solved: gamma transport (PTCrysp)}

A 511~keV annihilation photon entering the crystal undergoes a sequence of Compton scatters,
generally terminated by photoelectric absorption (or by escape). The PTCrysp engine
(vendored from PTCryspMC) transports the gamma with XCOM cross sections and Klein--Nishina
sampling and produces, per event, a list of energy deposits
\begin{equation}
  \mathcal{D} = \{ (\mathbf{r}_k,\, E_k,\, t_k) \},
  \qquad \sum_k E_k \le 511~\mathrm{keV},
\end{equation}
where $\mathbf{r}_k$ is the interaction point (electrons deposit locally at this scale),
$E_k$ the energy transferred, and $t_k$ the interaction time. Typically 1--3 deposits carry
almost all the energy. This list is the \emph{input} to CryspLight.

\subsection{What is new: optical transport}

Each deposit $k$ converts into scintillation photons:
\begin{equation}
  N_k \sim \mathrm{Poisson}\!\left( Y \cdot E_k \right),
\end{equation}
with $Y$ the light yield (Table~\ref{tab:crystals}). For a full-energy 511~keV event this
means $\sim 5\times10^4$ photons in cold \csi{} and $\sim 7\times10^3$ in \bgo{} at 195~K.
Each photon is emitted isotropically from $\mathbf{r}_k$, with a wavelength drawn from the
emission spectrum and an emission time
\begin{equation}
  t_{\gamma} = t_k + \Delta t, \qquad \Delta t \sim \mathrm{Exp}(\tau),
  \label{eq:emission-time}
\end{equation}
with $\tau = 800$~ns (\csi{}, 77~K) or $\tau = 1500$~ns (\bgo{}, 195~K). The photon then
propagates in straight lines inside the crystal, reflecting at the walls, until one of two
terminal fates:
\begin{enumerate}
  \item \textbf{Absorption}: in the bulk (exponential attenuation with length
    $\Lambda_{\rm abs}$) or at a wall (each reflector bounce succeeds only with probability
    $R < 1$).
  \item \textbf{Detection}: the photon crosses the back face into the SiPM plane and is
    detected with the photon-detection efficiency $\mathrm{PDE}(\lambda)$, recorded as a hit
    $(i, j, t_{\rm det})$ on SiPM $(i,j)$.
\end{enumerate}
The observables we must deliver per event are the per-SiPM photon counts (the ``light map,''
from which a transverse position is reconstructed), the photon arrival times (from which the
event time is estimated, e.g.\ by first-photon or few-photon statistics), and, optionally, the
full per-photon record for detector-response studies.

\subsection{Why the GPU}

The optical stage dominates the computational cost: $5\times10^4$ photons per event, each
undergoing tens of wall bounces, versus a handful of gamma interactions. A CPU Geant4 optical
simulation of such an event takes $\mathcal{O}(1)$~s; studies with $10^5$--$10^7$ events (light
maps vs.\ interaction depth, timing resolution vs.\ threshold, reflector comparisons) are
impractical without acceleration. But the photons never interact with each other and the
geometry is a single axis-aligned box: the problem maps perfectly onto a GPU, one thread per
photon. The target platform is Apple silicon (Metal), which is where this project's
development machines live.

%=====================================================================
\section{Optical properties of the materials}
\label{sec:properties}

This section collects the parameters the simulation needs, with typical literature values.
Every number here is a \emph{default to be overridden by measurement}: bulk attenuation in
particular varies between crystal boules, and reflectivities depend on wrap thickness and
mounting. The parameter set is deliberately identical in structure to a Geant4 material +
optical-surface definition, so that any value used in a Geant4 cross-check
(Sec.~\ref{sec:validation}) maps one-to-one onto ours.

\subsection{Crystals}
\label{sec:crystals}

\begin{table}[htb]
  \centering
  \caption{Crystal bulk and scintillation parameters used as defaults.
  $n$ is quoted at the emission peak. Yields and time constants for cold operation follow the
  CRYSP measurements~\cite{soleti2024} and standard references~\cite{mikhailik}.}
  \label{tab:crystals}
  \begin{tabular}{lcc}
    \toprule
     & \csi{} (pure, 77--100 K) & \bgo{} (195 K) \\
    \midrule
    Density (g/cm$^3$)                  & 4.51  & 7.13 \\
    Radiation length $\xzero$ (mm)      & 18.6  & 11.2 \\
    Depth $L_z = 2\xzero$ (mm)          & 37.2  & 22.4 \\
    Emission peak $\lambda_{\rm em}$ (nm) & $\sim$340 & $\sim$480 \\
    Emission range (nm)                 & 280--420 & 375--650 \\
    Refractive index $n(\lambda_{\rm em})$ & $\simeq 1.95$ & $\simeq 2.15$ \\
    Critical angle $\theta_c = \arcsin(1/n)$ & $30.9^\circ$ & $27.7^\circ$ \\
    Light yield $Y$ (photons/MeV)       & $\sim 1.0\times10^5$ & $\sim 1.5\times10^4$ \\
    Decay time $\tau$ (ns)              & 800   & 1500 \\
    Bulk attenuation $\Lambda_{\rm abs}$ (m) & $\gtrsim 0.5$ (measure) & $\sim 1$ (measure) \\
    \bottomrule
  \end{tabular}
\end{table}

Remarks:
\begin{itemize}
  \item \textbf{Refractive index.} Both crystals are strongly refractive, so total internal
    reflection (TIR) dominates the wall physics: any photon hitting a wall at incidence
    $\theta > \theta_c \simeq 28^\circ\text{--}31^\circ$ from the normal reflects with unit
    probability, regardless of the wrap. The wrap only matters for the minority cone
    $\theta < \theta_c$. The index should be tabulated vs.\ wavelength (a one-term Sellmeier
    fit is adequate over the emission band); dispersion matters mainly for the Fresnel
    transmission into the SiPM coupling layer.
  \item \textbf{Bulk attenuation.} At the emission wavelength, good crystals have
    $\Lambda_{\rm abs}$ of order a metre --- long compared with the crystal, but \emph{not}
    negligible: a trapped photon can travel many crystal lengths before escaping the TIR
    trap, so the effective path is tens of cm. $\Lambda_{\rm abs}(\lambda)$ should be a
    tabulated input, and is a prime candidate for in-situ calibration (tuning the simulated
    light map tail against data). \textbf{Rayleigh scattering} is implemented as a separate
    bulk process (length $\Lambda_{\rm R}$, phase function $\propto 1+\cos^2\theta$), and
    the absorption/scattering split matters more than their sum suggests: a measured
    ``attenuation length'' conflates the two, but scattering \emph{redirects} rather than
    kills --- it unlocks TIR-trapped photons (Sec.~\ref{sec:firstapp}) and restores
    ergodicity, so re-attributing part of a measured attenuation to scattering raises the
    predicted collection.
  \item \textbf{Time profile.} A single exponential per crystal for v1
    (Eq.~\ref{eq:emission-time}). The framework accepts a sum of exponentials with weights
    --- \bgo{} at lower temperatures develops a slow component, and cold \csi{} has a known
    fast/slow structure --- exactly like Geant4's
    \texttt{SCINTILLATIONCOMPONENT1/2} + \texttt{SCINTILLATIONTIMECONSTANT1/2}.
  \item \textbf{Emission spectrum.} Sampled once per photon at generation. In v1 a photon's
    wavelength enters through three lookups: $n(\lambda)$, $\Lambda_{\rm abs}(\lambda)$,
    $\mathrm{PDE}(\lambda)$, plus the ESR reflectivity edge (below). We precompute these at
    generation time and store per-photon scalars, so the kernel never interpolates tables.
\end{itemize}

\subsection{Reflectors: Vikuiti ESR and PTFE}
\label{sec:reflectors}

The two wrap materials behave very differently, and the difference interacts with the
crystal choice:

\textbf{3M Vikuiti ESR}~\cite{vikuiti} is a $\sim$65~$\mu$m all-polymer multilayer dielectric
mirror: a nearly ideal \emph{specular} reflector with $R > 98\%$ across 400--800~nm.
Crucially, being a polymeric interference stack, its reflectivity \emph{collapses below
$\sim$400~nm}~\cite{janecek2008}. For \bgo{} (480~nm emission) ESR is excellent. For pure cold
\csi{}, whose emission peaks at $\sim$340~nm, a large fraction of the spectrum falls below the
ESR edge, where the film is nearly transparent/absorbing: ESR is a poor default for cold
\csi{}, and this is precisely the kind of conclusion the simulation must be able to
demonstrate quantitatively. The wavelength-dependent $R_{\rm ESR}(\lambda)$ must therefore be
tabulated (measured curves in~\cite{janecek2008}), not taken as a constant.

\textbf{PTFE (Teflon)} is a \emph{diffuse} (to good approximation Lambertian) reflector whose
reflectivity saturates at $R \simeq 99\%$ for thicknesses $\gtrsim 0.5$--1~mm, and --- unlike
ESR --- it holds its reflectivity into the UV, down to $\sim$250~nm~\cite{janecek2008,
janecek2012}. It is the natural default for cold \csi{}.

\textbf{Baseline: thick PTFE on both crystals.} The initial production runs wrap both
crystals in 1--2~mm PTFE, i.e.\ at saturated, spectrally flat reflectivity; the ESR
comparison (and hybrid wraps) is a later study. This simplifies v0 twice over: a single
surface model (Lambertian back-painted) suffices, and with a flat reflector the residual
wavelength dependence of the transport reduces to bulk attenuation and the SiPM PDE. One
known approximation to record: a thick diffuse wrap returns some photons with a mm-scale
lateral offset (the light diffuses inside the PTFE before re-emerging). Standard practice
--- and v1 --- neglects it, re-injecting at the exit point; at 48~mm transverse size this
only slightly blurs the light-map edges.

Both are mounted with an air gap. The effective wall physics is therefore two-stage
(Sec.~\ref{sec:geant4}): first the crystal--air Fresnel/TIR decision at the polished
crystal surface, then, for the transmitted fraction, a reflection off the wrap with
probability $R_{\rm wrap}(\lambda)$ --- specular for ESR, Lambertian for PTFE --- after which
the photon re-enters the crystal.

\subsection{How Geant4 parameterizes this (and what we adopt)}
\label{sec:geant4}

In Geant4~\cite{geant4book}, the bulk is a \texttt{G4MaterialPropertiesTable} with
\texttt{RINDEX}, \texttt{ABSLENGTH}, \texttt{SCINTILLATIONCOMPONENT1},
\texttt{SCINTILLATIONYIELD}, \texttt{SCINTILLATIONTIMECONSTANT1},
\texttt{RESOLUTIONSCALE}; the surfaces are \texttt{G4OpticalSurface} objects under the
\textsc{unified} model~\cite{unified}. The \textsc{unified} finishes that describe a polished
crystal wrapped with an air-gapped reflector are exactly our two cases:
\begin{itemize}
  \item \texttt{polishedbackpainted}: polished crystal surface, air gap, \emph{specular}
    reflector of given \texttt{REFLECTIVITY} $\Rightarrow$ ESR;
  \item \texttt{groundbackpainted}: polished-or-ground crystal surface, air gap,
    \emph{Lambertian} reflector of given \texttt{REFLECTIVITY} $\Rightarrow$ PTFE.
\end{itemize}
In both, Geant4 first applies Snell/Fresnel (with TIR) at the crystal--air interface using
the two refractive indices, and only photons that escape into the gap see the reflector.
Two further \textsc{unified} ingredients are implemented, because they turn out to be
first-order for the light collection (Sec.~\ref{sec:firstapp}):
\begin{itemize}
  \item \textbf{Surface micro-roughness $\sigma_\alpha$} (the Gaussian width of the
    micro-facet normal distribution): each wall interaction samples a facet normal tilted
    by $\alpha \sim |\mathcal{N}(0,\sigma_\alpha)|$, applies Fresnel/TIR against the
    \emph{facet}, and reflects specularly about it (the specular lobe), with rejection of
    unilluminated facets and of reflections that would exit the surface. Measured polished
    scintillator surfaces have $\sigma_\alpha \approx 1.3^\circ$, ground ones
    $\sim 12^\circ$; $\sigma_\alpha$ is what allows TIR-trapped photons
    (Sec.~\ref{sec:firstapp}) to escape.
  \item \textbf{Front-painted finishes} (\texttt{polishedfrontpainted} /
    \texttt{groundfrontpainted}): the reflector in \emph{optical contact} --- no air gap,
    no Fresnel gate; every wall hit reaches the reflector (specular or Lambertian) with
    reflectivity $R$. This brackets the wrapping from the other side: real tape lies
    between the ideal gap and full contact.
\end{itemize}
More sophisticated measured-surface models exist --- the LUT Davis model~\cite{lutdavis}
ships lookup tables computed from AFM scans of real crystal surfaces (polished/etched/rough
$\times$ bare/ESR-air/ESR-glued/Teflon) and is available in Geant4 $\ge$10.3 and GATE; it is
the natural reference to tune our \textsc{unified} parameters against, alongside the Geant4
cross-check of Sec.~\ref{sec:validation}.

We adopt the \textsc{unified} physics (four finishes + $\sigma_\alpha$) as the
specification for the Metal kernel: it is well documented, widely validated, and directly
comparable with a Geant4 run using the same parameters.

\subsection{SiPM plane}
\label{sec:sipm}

The back face is coupled through optical grease ($n_g = 1.45$) to the SiPM windows (silicone
resin, $n \simeq 1.55$; the grease--window interface is nearly index-matched, reflectance
$< 0.2\%$, so the back face is modelled as a single crystal$\to$grease Fresnel surface).
A photon reaching the back face:
\begin{enumerate}
  \item undergoes the crystal$\to$grease Fresnel decision. Note that with
    $n_{\rm crystal} \simeq 2$ and $n_g = 1.45$ there is still a TIR cone:
    $\theta_c' = \arcsin(n_g/n) = 48.0^\circ$ (\csi{}) / $42.4^\circ$ (\bgo{}); photons
    beyond it reflect back into the crystal and keep bouncing --- they are \emph{not} lost;
  \item if transmitted, lands at $(x, y)$: mapped to SiPM indices
    $(i,j) = (\lfloor x/6~\mathrm{mm} \rfloor, \lfloor y/6~\mathrm{mm} \rfloor)$; the 6~mm
    pitch exactly tiles the face, so v1 takes the back plane as fully active, with a
    fill-factor test (dead border between packages) as a later config option;
  \item is detected with probability $\mathrm{PDE}(\lambda)$ (S14160-6050HS: $\simeq 0.5$ at
    the 450~nm peak, falling toward the UV; for 340~nm \csi{} light the PDE is substantially
    lower, another crystal--sensor coupling the simulation must expose). If not detected it
    is absorbed (v1; window reflection can be added later).
\end{enumerate}
\textbf{PDE is applied on the fly}, as the Bernoulli draw of step 3 inside the kernel, and
the outputs record \emph{detected} photons (photoelectrons) only. This choice is exact in
wavelength --- each photon carries its own $\mathrm{PDE}(\lambda)$, looked up at generation
--- and exact for first-photon statistics --- the earliest recorded time \emph{is} the
first photoelectron, with no downstream thinning bias. The price is that a production is
tied to one sensor's PDE curve; changing SiPM means re-running, which at the throughputs of
Sec.~\ref{sec:metal} is cheap. (Sensor effects beyond photon detection --- dark counts,
crosstalk, afterpulsing, pulse shape --- are downstream detector response, not simulated
here.)

The recorded hit is $(i, j, t_{\rm det})$ with
$t_{\rm det} = t_\gamma + \sum_{\rm segments} n\,\ell/c$, $c = 299.79$~mm/ns. Photon
propagation adds only ns-scale spreads, small against $\tau$, but essential for first-photon
timing studies.

%=====================================================================
\section{The transport algorithm}
\label{sec:algorithm}

\subsection{Photon state and geometry}

The geometry is a single axis-aligned box $[0,48]\times[0,48]\times[0,L_z]$~mm (crystal frame;
the SiPM plane at $z = L_z$). A photon is
\begin{equation}
  \gamma = (\mathbf{r},\, \hat{\mathbf{u}},\, t,\, w),
\end{equation}
position, unit direction, running time, and the per-photon wavelength-derived scalars
$w = (n, \Lambda_{\rm abs}, R_{\rm wrap}, \mathrm{PDE})$ precomputed at generation
(Sec.~\ref{sec:crystals}). No acceleration structures, no volume hierarchy: the distance to
the wall along $\hat{\mathbf{u}}$ is the closed-form axis-aligned slab test (3 divisions,
3 min/max).

\subsection{The bounce loop}

Each photon runs the same loop until it terminates or reaches a bounce cap
$B_{\max}$ (default 1000, never reached in practice with $R<1$):

\begin{enumerate}
  \item \textbf{Free flight.} Draw the absorption and Rayleigh distances
    $d_{\rm abs} = -\Lambda_{\rm abs} \ln \xi$, $d_{\rm R} = -\Lambda_{\rm R} \ln \xi'$;
    compute the wall distance $d_{\rm wall}$. If $d_{\rm abs}$ wins: bulk absorption,
    terminate (\textsc{absorbed-bulk}). If $d_{\rm R}$ wins: advance, redraw the direction
    from the $1+\cos^2\theta$ phase function, continue.
  \item \textbf{Advance.} $\mathbf{r} \mathrel{+}= d_{\rm wall}\hat{\mathbf{u}}$,
    $t \mathrel{+}= n\, d_{\rm wall}/c$. Identify the face and the incidence angle
    $\cos\theta = |\hat{\mathbf{u}}\cdot\hat{\mathbf{n}}|$.
  \item \textbf{Back face} ($z = L_z$): apply the crystal$\to$grease Fresnel decision
    (unpolarized Fresnel average, Eq.~\ref{eq:fresnel}, with TIR for
    $\sin\theta > n_g/n$). Reflected $\to$ flip $u_z$, continue. Transmitted $\to$
    SiPM-index lookup, fill-factor test, PDE Bernoulli: \textsc{detected}$(i,j,t)$ or
    \textsc{absorbed-sipm}.
  \item \textbf{Wrapped faces}, front-painted finish (contact): the reflector acts on
    every hit --- survive with probability $R_{\rm wrap}$ (specular or Lambertian
    re-injection), else \textsc{absorbed-wall}. No Fresnel gate.
  \item \textbf{Wrapped faces}, back-painted finish (air gap): two-stage \textsc{unified}
    physics:
    \begin{enumerate}
      \item Crystal--air Fresnel with $n_{\rm air}=1$, against the micro-facet normal when
        $\sigma_\alpha > 0$ (specular-lobe reflection about the facet, with the rejection
        rules of Sec.~\ref{sec:geant4}). TIR ($\sin\theta > 1/n$, the majority
        of hits): specular reflection, continue at zero cost. Otherwise
        reflect with the Fresnel probability or transmit into the gap.
      \item In the gap: survive with probability $R_{\rm wrap}$, else
        \textsc{absorbed-wall}. Surviving photons re-enter the crystal: for ESR, specular
        --- by symmetry of the two air crossings the photon re-enters with its transmitted
        direction mirrored, which reduces to a net specular flip of the original direction;
        for PTFE, draw a Lambertian direction about the inward normal
        ($u_\parallel$ from $\sqrt{\xi}$-weighted polar sampling) and refract it back
        through the surface (in v1 we adopt the standard simplification of re-injecting the
        Lambertian direction directly, as \textsc{unified} does for back-painted finishes;
        the gap geometry itself --- photon translation along the wrap --- is negligible for
        a $\lesssim 0.1$~mm gap).
    \end{enumerate}
\end{enumerate}

The unpolarized Fresnel reflectance between indices $n_1 \to n_2$ at incidence $\theta_i$
(refraction angle $\theta_t$ from Snell) is
\begin{equation}
  R_F = \frac{1}{2}\left[
    \left(\frac{n_1\cos\theta_i - n_2\cos\theta_t}{n_1\cos\theta_i + n_2\cos\theta_t}\right)^2
    +
    \left(\frac{n_1\cos\theta_t - n_2\cos\theta_i}{n_1\cos\theta_t + n_2\cos\theta_i}\right)^2
  \right].
  \label{eq:fresnel}
\end{equation}
Polarization is not tracked (standard for scintillator MC at this fidelity; Geant4 tracks it
but the effect on collection efficiency in a high-TIR box is small).

\subsection{Expected bounce statistics}

A photon inside the TIR cone of the back face ($\theta < \theta_c'$ w.r.t.\ $z$) exits after
$\mathcal{O}(1)$ bounces. The trapped majority random-walks until the wrap losses
($1 - R \simeq 1$--2\% per sub-critical wall hit) or bulk absorption kill it, or until a
diffuse (PTFE) bounce redirects it into the escape cone. Mean bounce counts are
$\mathcal{O}(10\text{--}100)$; the light collection efficiency at the SiPM plane is expected
in the 30--70\% range depending on wrap and crystal, with strong depth dependence --- these
distributions are primary outputs of the simulation, not inputs.

%=====================================================================
\section{GPU implementation on Metal}
\label{sec:metal}

\subsection{Language and stack}

We propose \textbf{Julia + Metal.jl}, writing the kernel with KernelAbstractions.jl so the
\emph{same} kernel source runs on the CPU (reference/debug path, thread-parallel) and on the
Metal backend, and ports to CUDA later if needed. This keeps CryspLight in the same language
as the vendored PTCrysp engine (Julia), with one build system and shared I/O. The alternative
--- a standalone Swift/C++ core with metal-cpp --- would give marginally more control over
kernel scheduling at the cost of a language boundary through the middle of the physics; we
take it only if Metal.jl proves limiting.

\subsection{Design constraints from Metal}

\begin{itemize}
  \item \textbf{Float32 only.} Metal GPUs have no usable Float64. This is comfortable here:
    positions in mm need $\sim$$10^{-4}$~mm precision (Float32 $\epsilon$ at 48~mm:
    $4\times10^{-6}$~mm); times are kept \emph{relative to the gamma interaction}, so even at
    $t \sim 10^4$~ns the Float32 resolution is $\sim$$10^{-3}$~ns, far below any physical
    jitter. Absolute event clocks stay on the CPU in Float64, exactly as PTCrysp already
    does.
  \item \textbf{No device-side RNG library.} We inline a counter-based RNG
    (Philox\,$4\times32$-10~\cite{philox}): stateless, keyed by (event id, photon id, draw
    counter), giving bit-reproducible results independent of thread scheduling --- important
    for regression tests and for resuming/replaying single photons on the CPU path.
  \item \textbf{Divergence.} The bounce loop has data-dependent length. With one thread per
    photon, a warp lives as long as its slowest photon. The bounce-count distribution has a
    modest tail (absorption caps it geometrically), so v1 accepts the divergence; if
    profiling shows it matters, the standard fix is a persistent-kernel/compaction scheme
    (terminate-and-refill from a global photon queue). This is an optimization, not a design
    change.
  \item \textbf{Memory layout.} Structure-of-arrays MtlArrays:
    $x, y, z, u_x, u_y, u_z, t, n, \Lambda, R, \mathrm{pde}$ (Float32) + photon id
    (UInt32). A batch is \emph{all photons of many events} concatenated ($10^6$--$10^7$
    photons per launch), with an event-offset table; per-photon state is
    $\sim$48~bytes, so even $10^7$ photons is $\sim$0.5~GB --- fine for unified memory.
\end{itemize}

\subsection{Kernel structure and outputs}

Generation (CPU or a trivial first kernel): expand deposits into photons --- Poisson yield,
isotropic direction, exponential emission time, wavelength draw, property lookups
($n$, $\Lambda_{\rm abs}$, $\mathrm{PDE}$; the wavelength itself is consumed here and never
stored). Transport (the kernel of Sec.~\ref{sec:algorithm}): one thread per photon, bounce
loop, terminal write.

\textbf{The standard output is fixed-size per event}, filled by device-side atomics,
recording detected photoelectrons (PDE already applied, Sec.~\ref{sec:sipm}):
\begin{enumerate}
  \item \textbf{Time-binned count map}: an $8\times8\times N_t$ array of photoelectron
    counts in \textbf{100~ns bins}, UInt32 on device (UInt16 on disk), window
    $\sim 5\tau$ plus an overflow bin: $N_t = 40$ (0--4~$\mu$s, \csi{}) or 80
    (0--8~$\mu$s, \bgo{}). One atomic add per detected photon;
    $\sim$5--10~kB per event.
  \item \textbf{First-photoelectron matrix}: an $8\times8$ Float32 array of the earliest
    detected-photon time per SiPM, initialized to $+\infty$ and filled by atomic-min.
    (Metal has no float atomic-min; for positive floats the IEEE-754 bit pattern is
    monotonic, so atomic-min on the value reinterpreted as UInt32 is exact.) Because PDE is
    applied in-kernel, this \emph{is} the first photoelectron --- no thinning bias. The
    100~ns histogram carries no sub-bin timing, so this matrix is the sole carrier of
    fast-timing information.
  \item \textbf{Summary counters}: generated / detected / absorbed-bulk / absorbed-wall /
    absorbed-sipm, per event.
\end{enumerate}
Two optional extensions, off by default: the \textbf{first-$k$} detected times per SiPM
($k \sim 8$--16, a fixed-size insertion sort in registers) for $k$-th--photoelectron trigger
emulation; and the \textbf{per-photon record} (status, SiPM index, detection time ---
8~bytes/photon) for validation samples, never productions.

\subsection{Throughput estimate}

Per bounce: one slab test, 1--3 RNG draws, one Fresnel evaluation --- order
$10^2$ FLOPs and no global-memory traffic beyond the initial state load (the loop runs in
registers). An M-class GPU sustains $\mathcal{O}(10)$ TFLOP/s FP32; even at 1\% arithmetic
efficiency that is $\gtrsim 10^9$ bounce-steps/s. A full-energy \csi{} event
($5\times10^4$ photons $\times$ $\sim$30 bounces $\simeq 1.5\times10^6$ steps) then costs
$\sim$1~ms of GPU time --- $10^3$ events/s, roughly three orders of magnitude over a
single-core Geant4 optical simulation, and \bgo{} is $7\times$ cheaper still. These are
order-of-magnitude figures to be confirmed by the first benchmark; they set the expectation
that $10^6$-event studies are minutes, not days.

%=====================================================================
\section{Interfaces}
\label{sec:interfaces}

\textbf{Input} --- the PTCrysp deposit list per event (HDF5): event id, deposit position
(mm, crystal frame), energy (keV), time (ns). Plus a TOML config: crystal (selects
Table~\ref{tab:crystals} defaults), wrap per face group (ESR / PTFE / none), coupling and
SiPM parameters, transport options, batch size, seed.

\textbf{Output} --- per event (HDF5): the $8\times8\times N_t$ photoelectron count map
(100~ns bins), the $8\times8$ first-photoelectron time matrix, the summary counters, and
optionally the first-$k$ times and the per-photon record. PDE is already applied, so these
are detected quantities for the configured sensor; the sensor name and PDE-curve version are
stamped, with config and seed, as root attributes following the PTCrysp provenance
convention.

%=====================================================================
\section{Validation plan}
\label{sec:validation}

\begin{enumerate}
  \item \textbf{Analytic unit tests} (CPU path, then bit-compared with Metal):
    Fresnel/Snell values against closed forms; TIR cone fractions ($f = 1 - \cos\theta_c$
    per face for an isotropic source) with reflectivities set to 1 and absorption off;
    exponential pathlength distribution against $\Lambda_{\rm abs}$; energy conservation
    (generated = detected + absorbed, exactly); Lambertian angular distribution of PTFE
    bounces.
  \item \textbf{Limiting configurations}: $R=1$, $\Lambda_{\rm abs}=\infty$, full-face
    perfect detector $\Rightarrow$ every photon must be detected; a black box
    $\Rightarrow$ only the direct cone arrives, with the closed-form solid-angle count and
    arrival-time distribution.
  \item \textbf{Geant4 cross-check} (the decisive one): a standalone Geant4 application with
    the identical box, the \textsc{unified} back-painted surfaces of
    Sec.~\ref{sec:geant4}, and the same property tables. Compare, vs.\ interaction depth
    and transverse position: collection efficiency, light-map shapes (row/column profiles),
    arrival-time spectra, wall-loss fractions. Agreement at the few-percent level validates
    the kernel physics; residual differences localize to a specific surface parameter by
    construction, since the parameter sets are isomorphic.
  \item \textbf{Data anchors}: the CRYSP cold-\csi{} measurements~\cite{soleti2024} (light
    collection, energy resolution, timing with the same SiPM family) provide the
    experimental anchor for the \csi{} + PTFE configuration; the ESR-vs-PTFE UV behaviour
    checks against the measured reflectances of~\cite{janecek2008}.
\end{enumerate}

\subsection{First application: centred isotropic source}
\label{sec:firstapp}

The first end-to-end application (built, with the v0 CPU engine) is an isotropic point
source at the crystal centre with $\mathrm{PDE}=1$ and \emph{scalar} optics --- each
material reduced to single numbers, fixed at the emission wavelength: \csi{}
($\lambda=350$~nm, $n=1.95$, $\tau=800$~ns, $Y=10^5$/MeV), \bgo{} ($\lambda=550$~nm,
$n=2.15$, $\tau=1500$~ns, $Y=1.5\times10^4$/MeV), grease $n_g=1.45$, Teflon
$R=0.99$ (Lambertian), ESR $R=0.985$ (specular). Three cases: \csi{}+Teflon,
\bgo{}+Teflon, and \bgo{}+ESR --- the last exercises the specular surface branch with a
crystal for which ESR is legitimate. Outputs are the standard matrices
(Sec.~\ref{sec:metal}).

Two exact anchors follow from the fact that specular interactions (including TIR)
preserve the direction components $(|u_x|,|u_y|,|u_z|)$:
\begin{itemize}
  \item \textbf{Idealized specular limit} ($R=1$, $\Lambda_{\rm abs}=\infty$,
    $\mathrm{PDE}=1$, ESR): the detected fraction is \emph{exactly} the grease escape
    cone, $1-\cos\theta_c'$ (0.331 for \csi{}, 0.261 for \bgo{}); everything else bounces
    forever.
  \item \textbf{Trapped direction classes}: with a polished crystal and an air gap, a
    photon super-critical on \emph{all} faces never reaches the wrap at all --- the
    Lambertian randomization only acts on photons inside a crystal--air Fresnel cone. In
    the idealized Teflon limit these photons are never detected; with physical parameters
    they die by bulk absorption. Consequently the light collection of an ideally polished,
    air-gapped crystal is dominated by the interplay of trapping and
    $\Lambda_{\rm abs}$, and the surface roughness $\sigma_\alpha$ (which unlocks the
    trapped classes) is a first-order parameter for the data comparison, not a refinement.
\end{itemize}
Both anchors are enforced in the test suite, together with exact photon conservation
across the terminal statuses, the Fresnel closed forms, the sampling distributions, and two
further exact limits: \emph{Rayleigh restores ergodicity} (pure scattering, perfect wrap,
$\mathrm{PDE}=1$ $\Rightarrow$ every photon is detected --- no direction class survives
redirection), and the contact-Lambertian $R=1$ limit detects everything (no Fresnel gate,
no trapping).

\paragraph{Sensitivity of the light collection to the realism ingredients.}
Table~\ref{tab:surfscan} (generated by \texttt{scripts/study\_surface.jl};
\csi{}+Teflon, centred source, $\mathrm{PDE}=1$, $\delta$ emission) quantifies each
ingredient. The lessons: (i) the ideally polished air-gapped crystal is the pessimistic
extreme --- its collection is capped near the escape cone; (ii) $\sigma_\alpha$ is
first-order: the measured-polish $1.3^\circ$ already moves the collection by 6 points, a
ground surface by nearly 30; (iii) Rayleigh scattering acts in the same direction
(redirection unlocks trapped light); (iv) $\Lambda_{\rm abs}$ and $\sigma_\alpha$
\emph{interact}: a long $\Lambda_{\rm abs}$ alone barely helps the polished crystal (the
trapped classes still die in the bulk, just later), but combined with even $1.3^\circ$ of
roughness it yields the highest collection of the scan --- the unlocking random walk needs
path length to work with. Consequently the pair
$(\Lambda_{\rm abs}, \sigma_\alpha)$ must be calibrated \emph{jointly} against data
(collection + light-map shape + time tail), not one at a time.

\begin{table}[htb]
  \centering
  \caption{\csi{}+Teflon light collection vs.\ surface/bulk model
  ($10^6$ photons, centred point source, $\mathrm{PDE}=1$, $R_{\rm wrap}=0.99$,
  $\Lambda_{\rm abs}=500$~mm unless stated).}
  \label{tab:surfscan}
  \begin{tabular}{lccc}
    \toprule
    Configuration & detected & bulk-absorbed & wall-absorbed \\
    \midrule
    air gap, polished ($\sigma_\alpha=0$, v0)        & 34.1\% & 65.5\% & 0.5\% \\
    air gap, $\sigma_\alpha = 1.3^\circ$             & 39.8\% & 59.7\% & 0.6\% \\
    air gap, $\sigma_\alpha = 6^\circ$               & 55.8\% & 43.3\% & 0.9\% \\
    air gap, $\sigma_\alpha = 12^\circ$ (ground)     & 62.9\% & 36.1\% & 1.0\% \\
    contact (front-painted) Lambertian               & 63.9\% & 31.8\% & 4.3\% \\
    air gap, polished + Rayleigh 300~mm              & 50.4\% & 48.9\% & 0.7\% \\
    air gap, $1.3^\circ$ + Rayleigh 300~mm           & 52.2\% & 47.0\% & 0.7\% \\
    air gap, polished, $\Lambda_{\rm abs}=5$~m       & 38.5\% & 61.0\% & 0.5\% \\
    air gap, $1.3^\circ$, $\Lambda_{\rm abs}=5$~m    & 69.9\% & 29.1\% & 1.1\% \\
    \bottomrule
  \end{tabular}
\end{table}

The baseline configs adopt the measured-polish $\sigma_\alpha = 1.3^\circ$ with the air
gap; the contact mode and $\Lambda_{\rm R}$ stay available as calibration freedom.

%=====================================================================
\section{Roadmap}
\label{sec:roadmap}

\begin{enumerate}
  \item \textbf{v0 --- CPU reference} (Julia, KernelAbstractions CPU backend): full physics
    of Sec.~\ref{sec:algorithm}, per-photon records, unit tests 1--2 above. Small but
    complete.
  \item \textbf{v1 --- Metal port}: same kernel on MtlArrays; bit-level agreement with v0 on
    identical Philox streams; first throughput benchmark.
  \item \textbf{v2 --- PTCrysp integration}: vendor the gamma engine, read its deposits,
    produce the per-event light maps + times end to end; Geant4 cross-check campaign.
  \item \textbf{v3 --- physics studies}: light map vs.\ depth-of-interaction (input to
    position reconstruction), first-photon timing vs.\ yield and $\tau$, ESR vs.\ PTFE
    vs.\ hybrid wraps for both crystals, sensitivity to $\Lambda_{\rm abs}$ and fill
    factor.
\end{enumerate}

%=====================================================================
\begin{thebibliography}{9}

\bibitem{soleti2024}
S.~R. Soleti \emph{et al.},
``Cryogenic CsI as a potential PET material,''
arXiv:2406.13598 (2024).

\bibitem{mikhailik}
V.~B. Mikhailik and H.~Kraus,
``Scintillators for cryogenic applications: state-of-art,''
J.~Phys.\ Stud.\ \textbf{14} (2010) 4201; and references therein for low-temperature
CsI and BGO yields and decay times.

\bibitem{janecek2008}
M.~Janecek and W.~W. Moses,
``Optical reflectance measurements for commonly used reflectors,''
IEEE Trans.\ Nucl.\ Sci.\ \textbf{55} (2008) 2432.

\bibitem{janecek2012}
M.~Janecek,
``Reflectivity spectra for commonly used reflectors,''
IEEE Trans.\ Nucl.\ Sci.\ \textbf{59} (2012) 490.

\bibitem{vikuiti}
3M Optical Systems Division,
``Vikuiti\texttrademark\ Enhanced Specular Reflector (ESR),'' product datasheet.

\bibitem{geant4book}
Geant4 Collaboration,
``Geant4 Book For Application Developers,'' ch.\ ``Physics Processes / Optical Photon
Processes'' (Rel.\ 11.x).

\bibitem{unified}
A.~Levin and C.~Moisan,
``A more physical approach to model the surface treatment of scintillation counters and its
implementation into DETECT,''
IEEE Nucl.\ Sci.\ Symp.\ Conf.\ Rec.\ (1996) 702.

\bibitem{lutdavis}
E.~Roncali \emph{et al.},
``An integrated model of scintillator-reflector properties for advanced simulations of
optical transport,''
Phys.\ Med.\ Biol.\ \textbf{62} (2017) 4811 (the LUT Davis model).

\bibitem{hamamatsu}
Hamamatsu Photonics,
``MPPC S14160/S14161 series,'' product datasheet.

\bibitem{philox}
J.~K. Salmon \emph{et al.},
``Parallel random numbers: as easy as 1, 2, 3,''
Proc.\ SC'11 (2011).

\end{thebibliography}

\end{document}
